\section{Conclusion} 
\label{sec:conclusion}


In this work, we quantified the behavioral and performance cost of consuming digital layout-based media when visual access is constrained. This challenge affects both readers on small-screen devices and people with low vision.

Our multidimensional study provides empirical evidence 
and quantifies the impact of different newspaper interaction modalities. We found that while manual magnification (\condone) is the standard, it imposes a navigational overhead that disrupts the reader's canonical strategy. In contrast, \condtwo acts as a modality that restores behavioral integrity.
 Compared to gesture-based magnification, \condtwo enabled faster performance and a natural reading path (closely resembling the reference strategy), a stronger mental representation of document structure, and a significantly lower perceived workload.
 
These results indicate that magnification alone may be insufficient to support effective reading and consistent quality of experience in layout-based documents. The standard paradigm (\condone) can make navigating a page more challenging, not due to a lack of features, but by potentially disrupting spatial context and forcing a shift from natural content consumption. By contrast, re-layouting documents into large-print editions (\condtwo) preserves the spatial qualities of the layout (such as entry points) and supports efficient, predictable reading strategies.



Our findings offer important implications: designing \replacedimx{digital media consumption tools}{for accessibility} should go beyond interactive magnification tools toward adaptive layouts that combine font scaling with structural re-flow. Our findings confirm that the  effort required to generate \condtwo is justified \addedimx{not just} by the improvements in reader performance, \addedimx{but also by the measurable restoration of behavioral integrity} and satisfaction across diverse populations and devices, supporting efficient and inclusive access to newspapers, magazines, and other visually complex digital content. While our study focused on text-only layout-based documents, future work should extend these methods to richer layouts, incorporating graphical elements such as images and advertisements, and should include careful empirical studies to evaluate additional types of entry points.